\section{Where meaning comes from}\label{sec:meaning}
A ledger needs a meaning $W$. There are two ways to obtain one. Someone can declare it, or it can be derived from what the message is for.

\subsection{Declared meaning}
In many systems $W$ is attached to $X$ by a process outside the source. The genetic code is a physical process. A lexicographer's sense inventory is a human one: WordNet assigns each word occurrence a synset \citep{miller1995,fellbaum1998}. The behaviour of a program on a test suite is an operational one.
In each case the ledger of \cref{sec:ledger} applies as it stands. What it cannot say is whether the declared $W$ is the right one for a given use. That is the question of the next two subsections.

\subsection{Meaning for a purpose}
Wittgenstein wrote that the meaning of a word is its use in the language \citep{wittgenstein1953}, and Bateson called information a difference which makes a difference \citep{bateson1972}. Both point to consequences. We take them literally.

\begin{definition}[Purpose, belief, meaning for a purpose]\label{def:purpose}
Let $(X,Y)$ be jointly distributed with finite ranges; $Y$ is the \emph{purpose}, whatever the receiver will act on or learn. For $x$ in the support of $X$ the \emph{belief} is $b(x)=\Pr(Y\in\cdot\given X{=}x)$, a distribution on $\mathcal Y$.
The \emph{meaning of $X$ for purpose $Y$} is the random element $M_Y=b(X)$; equivalently, the partition of the support of $X$ into the level sets of $b$.
\end{definition}
Two messages have the same meaning for $Y$ exactly when they leave the receiver with the same belief about $Y$. This depends only on $P_{XY}$; no annotator is involved.

\begin{theorem}[Meaning for a purpose]\label{thm:derived}
Let $M=M_Y$.
\begin{enumerate}[label=\textup{(\alph*)}]
\item \emph{Sufficiency.} $Y-M-X$ is a Markov chain. Hence $I(X;Y\given M)=0$ and $I(M;Y)=I(X;Y)$.
\item \emph{Minimality.} If $T=f(X)$ is any function with $I(T;Y)=I(X;Y)$, then $M$ is a function of $T$ almost surely. So $M$ is the coarsest function of $X$ that loses no information about $Y$; it is unique up to relabelling.
\item \emph{Ledger.} For the pair $(M,X)$ the doubt is zero and
\begin{equation}
E=H(M)+H(X\given M)=\underbrace{I(X;Y)}_{\text{consequence}}+\underbrace{H(M\given Y)}_{\text{hedge}}+\underbrace{H(X\given M)}_{\text{latitude}} .
\end{equation}
\item \emph{Monotone in the purpose.} If $Y'=h(Y)$ for a function $h$, then $M_{Y'}$ is a function of $M_Y$; so $H(M_{Y'})\le H(M_Y)$ and $H(X\given M_{Y'})\ge H(X\given M_Y)$.
\item \emph{Extremes.} If $Y$ is constant, $M$ is constant and all of $E$ is latitude. If $Y=X$, then $M$ is in one-to-one correspondence with $X$ and the latitude is zero.
\end{enumerate}
\end{theorem}
\begin{proof}
(a) Let $m$ be a value of $M$ and $x$ any message with $b(x)=m$. Then $\Pr(y\given x,m)=\Pr(y\given x)=m(y)$, and $\Pr(y\given m)=\sum_{x:\,b(x)=m}\Pr(x\given m)\,b(x)(y)=m(y)$. So $\Pr(y\given x,m)=\Pr(y\given m)$, that is $Y\perp X\mid M$. Since $M$ is a function of $X$, the chain rule gives $I(X;Y)=I(X,M;Y)=I(M;Y)+I(X;Y\given M)=I(M;Y)$.
(b) Let $T=f(X)$ with $I(T;Y)=I(X;Y)$. Because $T$ is a function of $X$, $I(X;Y\given T)=I(X,T;Y)-I(T;Y)=I(X;Y)-I(T;Y)=0$, so $Y\perp X\mid T$. Hence for every $t$ with positive probability and all $x,x'$ in $f^{-1}(t)$ of positive probability, $\Pr(y\given x)=\Pr(y\given t)=\Pr(y\given x')$. Thus $b$ is constant on each fibre of $f$, that is $M=g(T)$.
(c) $M$ is a function of $X$, so $H(M\given X)=0$; the chain rule gives $H(X)=H(X,M)=H(M)+H(X\given M)$. By (a), $H(M)=I(M;Y)+H(M\given Y)=I(X;Y)+H(M\given Y)$.
(d) The belief about $Y'$ is the image of the belief about $Y$ under $h$, so $b(x)=b(x')$ implies $b'(x)=b'(x')$. A function of a variable has no more entropy, and conditioning on more leaves no more.
(e) If $Y$ is constant, all beliefs coincide. If $Y=X$, then $b(x)$ is the point mass at $x$, which determines $x$.
\end{proof}
The consequence term $I(X;Y)$ is the information about $Y$ that the message carries, and it is the only one of the three that grows when the message becomes more informative. The hedge $H(M\given Y)$ is the entropy of the receiver's confidence not explained by the truth: a forecaster who says ``30\,\%'' on some days and ``70\,\%'' on others has content $1$ bit in $M$ but conveys less than one bit about rain. Content, in this sense, counts distinct beliefs; it is not the same as information. The latitude is what remains: variation of the form to which the purpose is blind.

\subsection{Adequacy and leakage}
A declared meaning $W$ is adequate for a purpose when it screens the form off from the consequence.
\begin{definition}[Adequacy and leakage]\label{def:leak}
Let $W$ be jointly distributed with $(X,Y)$. The \emph{leakage} of $W$ about the purpose $Y$ is $\Lk_Y(W)=I(X;Y\given W)$. The meaning $W$ is \emph{adequate} for $Y$ if $\Lk_Y(W)=0$.
\end{definition}
\begin{theorem}[Leakage]\label{thm:leak}
\begin{enumerate}[label=\textup{(\alph*)}]
\item $\Lk_Y(W)\ge0$, with equality iff $Y\perp X\mid W$.
\item If $W=f(X)$, then $\Lk_Y(W)=I(X;Y)-I(W;Y)$, and $W$ is adequate for $Y$ iff $M_Y$ is a function of $W$, that is, iff the partition induced by $f$ refines the partition induced by $b$. In general $\Lk_Y(W)=I(X;Y)-I(W;Y)+I(W;Y\given X)$.
\item \emph{Abstraction leaks.} If $W_1=f(X)$ and $W_2=g(W_1)$, then $\Lk_Y(W_2)=\Lk_Y(W_1)+I(W_1;Y\given W_2)$: each step of abstraction leaks exactly the relevance it discards.
\item \emph{No universal meaning.} If $W=f(X)$ is adequate for every purpose $Y$, then $f$ is injective on the support of $X$.
\end{enumerate}
\end{theorem}
\begin{proof}
(a) is the non-negativity of conditional mutual information. (b) For $W=f(X)$ the chain rule gives $I(X;Y)=I(X,W;Y)=I(W;Y)+I(X;Y\given W)$. In general $I(X,W;Y)$ expands in both orders, \[ I(W;Y)+I(X;Y\given W)=I(X;Y)+I(W;Y\given X), \] which is the last formula. If $I(W;Y)=I(X;Y)$ then \cref{thm:derived}(b) gives $M=g(W)$. Conversely, if $M=g(W)$ then $I(W;Y)\ge I(M;Y)=I(X;Y)$ by data processing, and $I(W;Y)\le I(X;Y)$ because $W$ is a function of $X$.
(c) By (b), $\Lk_Y(W_2)-\Lk_Y(W_1)=I(W_1;Y)-I(W_2;Y)$, and $I(W_1;Y)=I(W_1,W_2;Y)=I(W_2;Y)+I(W_1;Y\given W_2)$ since $W_2$ is a function of $W_1$.
(d) Take $Y=X$. Then $\Lk_X(W)=I(X;X\given W)=H(X\given W)$, which vanishes iff $X$ is a function of $W=f(X)$, iff $f$ is injective on the support.
\end{proof}
Part (d) is the philosophical content of the section. There is no purpose-free meaning: a partition that discards anything discards it for some purpose. The trivial partition is adequate for the purpose ``nothing'', and the identity is adequate for the purpose ``the exact words''. Every meaning in between is a statement about what the receiver cares about.

The leakage has an everyday reading. Linguists call the part of a word's effect that its denotation does not fix its connotation, register or style. Leakage measures it in bits, for a stated purpose: how much a hearer learns about the consequence from the choice among synonyms.

\subsection{Depth}
Meanings come in levels. A codon has an amino acid, which has a chemical class; a word occurrence has a sense, which has a semantic field.
\begin{proposition}[Depth ledger]\label{prop:depth}
Let $W_1=f_1(X)$, $W_2=f_2(W_1)$, \dots, $W_L=f_L(W_{L-1})$ be successive functions, and $W_0=X$. Then
\begin{equation}
H(X)=H(W_L)+\sum_{k=1}^{L}H(W_{k-1}\given W_k),\qquad
I(X;Y)=I(W_L;Y)+\sum_{k=1}^{L}I(W_{k-1};Y\given W_k).
\end{equation}
The term with $k=1$, the first of each sum, is the latitude $H(X\given W_1)$ and the leakage $\Lk_Y(W_1)$, respectively: the contribution of the finest level.
\end{proposition}
\begin{proof}
Since $W_k$ is a function of $W_{k-1}$, the chain rule gives
\[ H(W_{k-1})=H(W_k)+H(W_{k-1}\given W_k),\qquad I(W_{k-1};Y)=I(W_k;Y)+I(W_{k-1};Y\given W_k). \]
Sum over $k$ and telescope.
\end{proof}
\cref{fig:depth} shows the two ledgers side by side. They are two decompositions of two different quantities, the entropy of the message and the information it carries about a purpose, over the same levels. The first says where the choices are made; the second says where the consequences live. In \cref{sec:measurements} the two profiles differ sharply: the entropy of a codon is spread over all three levels, while the information it carries about which genome it came from sits almost entirely at the finest level, the synonymous choice.
\begin{figure}[tbp]
\centering
\begin{tikzpicture}[x=1cm,y=1cm,font=\footnotesize\sffamily]
  \definecolor{lvlA}{HTML}{2B5C8A}\definecolor{lvlB}{HTML}{6E9BC5}
  \node[anchor=west,font=\footnotesize\sffamily\bfseries] at (0.2,1.55) {Two ledgers on the same three levels of meaning\enspace\normalfont\scriptsize(not to scale)};
  % column headers: the levels
  \foreach \x/\t in {4.2/{level 2}, 6.8/{level 1}, 8.95/{level 0}}{
    \node[anchor=south,align=center,text=mute,font=\scriptsize\sffamily] at (\x,0.42) {\t};
  }
  % row 1: entropy of the form
  \node[anchor=east,align=right,text=mute] at (2.55,0) {entropy\\of the form};
  \fill[lvlA] (2.8,-0.3) rectangle (5.6,0.3);\fill[lvlB] (5.6,-0.3) rectangle (8.0,0.3);\fill[form] (8.0,-0.3) rectangle (9.9,0.3);
  \draw[white,line width=0.8pt] (5.6,-0.3) -- (5.6,0.3) (8.0,-0.3) -- (8.0,0.3);
  \node[text=white] at (4.2,0) {$H(W_2)$};\node[text=white] at (6.8,0) {$H(W_1\!\mid\! W_2)$};\node[text=ink] at (8.95,0) {$H(X\!\mid\! W_1)$};
  \node[anchor=west,text=mute] at (10.05,0) {$=H(X)$};
  % row 2: information about a purpose Y
  \node[anchor=east,align=right,text=mute] at (2.55,-1.5) {information\\about $Y$};
  \fill[lvlA] (2.8,-1.8) rectangle (4.4,-1.2);\fill[lvlB] (4.4,-1.8) rectangle (6.7,-1.2);\fill[leak] (6.7,-1.8) rectangle (9.9,-1.2);
  \draw[white,line width=0.8pt] (4.4,-1.8) -- (4.4,-1.2) (6.7,-1.8) -- (6.7,-1.2);
  \node[text=white,font=\scriptsize\sffamily] at (3.6,-1.5) {$I(W_2;Y)$};\node[text=white,font=\scriptsize\sffamily] at (5.55,-1.5) {$I(W_1;Y\!\mid\! W_2)$};
  \node[text=white,align=center,font=\scriptsize\sffamily] at (8.3,-1.5) {leakage $\Lk_Y(W_1)$\\$=I(X;Y\!\mid\! W_1)$};
  \node[anchor=west,text=mute] at (10.05,-1.5) {$=I(X;Y)$};
\end{tikzpicture}
\caption{Depth. If $W_1=f(X)$ and $W_2=g(W_1)$ are successively coarser meanings, the entropy of the form is a sum with one term per level (\cref{prop:depth}); the information the form carries about a purpose $Y$ splits over the same levels, and its finest-level term, drawn at the right, is the leakage of $W_1$ (the $k=1$ term of the sum). The two bars have different scales.}
\label{fig:depth}
\end{figure}

